\documentclass[11pt,reqno]{amsart}

\usepackage{amsmath,amssymb,amsthm}
\usepackage[margin=1.1in]{geometry}
\usepackage{microtype}
\usepackage[colorlinks=true,linkcolor=blue,urlcolor=blue]{hyperref}

\newtheorem{theorem}{Theorem}
\theoremstyle{remark}
\newtheorem{remark}[theorem]{Remark}

\newcommand{\Z}{\mathbb Z}
\newcommand{\Fp}{\mathbb F_p}
\newcommand{\Zloc}{\mathbb Z_{(p)}}
\DeclareMathOperator{\Num}{Num}

\title{A Prime-Power Congruence for an Odd-Exponent Binomial Family}
\author{}
\date{}

\begin{document}

\begin{abstract}
For a nonnegative integer $m$, define
$$
 u_m(n)=\Num\left(
 \sum_{k=1}^{n}\frac{1}{k^{2m+1}}
 \binom{n}{k}^{2}\binom{n+k}{k}^{2}
 \right),
$$
where the numerator is taken after reducing the rational number to
lowest terms.  The conjecture treated here originates in the record for
\href{https://oeis.org/A357513}{OEIS A357513}, which states this
parameterized form.  We determine exactly when $p^4$ divides
$u_m(p-1)$.
For every prime $p$,
$$
 p^4\mid u_m(p-1)
 \quad\Longleftrightarrow\quad
 p-1\nmid 2m+4\ \text{ or }\ p\mid 2m+7.
$$
Consequently, for each fixed $m$ there are only finitely many
exceptional primes; every such prime is at most $2m+5$.
\end{abstract}

\maketitle

\section{Statement}

The summation begins at $k=1$.  A version beginning at $k=0$ is not
defined because its first summand contains $1/0^{2m+1}$.

For a fixed nonnegative integer $m$, put
$$
 \mathcal E_m=
 \left\{p\ \text{prime}:
 p-1\mid 2m+4\ \text{and}\ p\nmid 2m+7\right\}.
$$

\begin{theorem}\label{thm:main}
Let $m\geq0$ be an integer and let $p$ be prime.  Then
$$
 p^4\mid u_m(p-1)
 \quad\Longleftrightarrow\quad
 p\notin\mathcal E_m.
$$
Moreover, $\mathcal E_m$ is finite and is contained in the primes at
most $2m+5$.  If $p\in\mathcal E_m$, then
$$
 v_p\bigl(u_m(p-1)\bigr)=
 \begin{cases}
  2,&p=2,\\
  3,&p\ \text{is odd}.
 \end{cases}
$$
\end{theorem}

\section{Common denominators}

Fix $m$, and write
$$
 d=2m+1.
$$
For a prime $p$, let
$$
 S_{m,p}=\sum_{k=1}^{p-1}\frac{1}{k^d}
 \binom{p-1}{k}^{2}\binom{p-1+k}{k}^{2}.
$$
When $p$ is odd, we work in the local ring
$$
 \Zloc=\left\{\frac{a}{b}:a,b\in\Z,\ p\nmid b\right\}.
$$

The reduced denominator of $S_{m,p}$ is coprime to $p$.  Indeed, set
$$
 D_m=((p-1)!)^d
$$
and
$$
 N_m=\sum_{k=1}^{p-1}
 \left(\binom{p-1}{k}\binom{p-1+k}{k}\right)^2
 \frac{D_m}{k^d}.
$$
Since $k\mid(p-1)!$ for $1\leq k\leq p-1$, every $D_m/k^d$ is an
integer.  Thus $N_m\in\Z$, $S_{m,p}=N_m/D_m$, and $p\nmid D_m$.

Suppose that $S_{m,p}=A/B$ in lowest terms.  Cross-multiplication gives
$AD_m=N_mB$.  Since $\gcd(A,B)=1$, Euclid's lemma gives $B\mid D_m$.
In particular, $p\nmid B$.  Thus both $B$ and $D_m$ are $p$-adic
units, so for every $r\geq0$,
\begin{equation}\label{eq:num-equiv}
 p^r\mid A
 \quad\Longleftrightarrow\quad
 S_{m,p}\equiv0\pmod{p^r}.
\end{equation}
In particular, a congruence for $S_{m,p}$ transfers directly to the
reduced numerator $u_m(p-1)$.

\section{The modular reduction}

For $1\leq k\leq p-1$, the exact identities
\begin{align*}
 \binom{p-1}{k}
 &=(-1)^k\prod_{j=1}^{k}\left(1-\frac pj\right),\\
 \binom{p-1+k}{k}
 &=\frac pk\prod_{j=1}^{k-1}\left(1+\frac pj\right)
\end{align*}
give
\begin{equation}\label{eq:binom-product}
 \binom{p-1}{k}\binom{p-1+k}{k}
 =(-1)^k\frac pk\left(1-\frac pk\right)
 \prod_{j=1}^{k-1}\left(1-\frac{p^2}{j^2}\right).
\end{equation}
After squaring \eqref{eq:binom-product} and dividing by $k^d$, the
$k$th summand of $S_{m,p}$ is
$$
 \frac{p^2}{k^{d+2}}\left(1-\frac pk\right)^2
 \prod_{j=1}^{k-1}\left(1-\frac{p^2}{j^2}\right)^2.
$$
All denominators are units in $\Zloc$, and
$$
 \prod_{j=1}^{k-1}\left(1-\frac{p^2}{j^2}\right)^2
 \equiv1\pmod{p^2}.
$$
The displayed summand already contains $p^2$.  Hence, modulo $p^4$, it
is congruent to
$$
 \frac{p^2}{k^{d+2}}-\frac{2p^3}{k^{d+3}}.
$$

Put
$$
 r=d+2=2m+3,\qquad e=d+3=2m+4
$$
and define
$$
 H_s=\sum_{k=1}^{p-1}\frac1{k^s}\in\Zloc.
$$
Summing the preceding congruence gives
\begin{equation}\label{eq:first-reduction}
 S_{m,p}\equiv p^2H_r-2p^3H_e\pmod{p^4}.
\end{equation}

The exponent $r$ is odd.  Since $k\mapsto p-k$ permutes
$1,\ldots,p-1$,
\begin{align*}
 \frac1{(p-k)^r}
 &=-\frac1{k^r}\left(1-\frac pk\right)^{-r}\\
 &\equiv-\frac1{k^r}-\frac{rp}{k^{r+1}}\pmod{p^2}.
\end{align*}
Because $r+1=e$, summing over $k$ yields
\begin{equation}\label{eq:harmonic-pairing}
 2H_r+rpH_e\equiv0\pmod{p^2}.
\end{equation}
Multiply \eqref{eq:harmonic-pairing} by $p^2$ and double
\eqref{eq:first-reduction}.  Since $r+4=2m+7$, we obtain
\begin{equation}\label{eq:key}
 2S_{m,p}\equiv-(2m+7)p^3H_e\pmod{p^4}.
\end{equation}

\section{The finite-field obstruction}

Reduction modulo $p$ and the fact that inversion permutes
$\Fp^\times$ give
$$
 H_e\equiv\sum_{x\in\Fp^\times}x^{-e}
 =\sum_{x\in\Fp^\times}x^e\pmod p.
$$
The standard power-sum evaluation in the cyclic group
$\Fp^\times$ is
\begin{equation}\label{eq:power-sum}
 H_e\equiv
 \begin{cases}
  0\pmod p,&p-1\nmid e,\\
  -1\pmod p,&p-1\mid e.
 \end{cases}
\end{equation}

Assume first that $p$ is odd.  Then $2$ is a unit in $\Zloc$.  If
$p-1\nmid e$, equations \eqref{eq:key} and \eqref{eq:power-sum} give
$S_{m,p}\equiv0\pmod{p^4}$.  If $p-1\mid e$, those equations instead
give
$$
 2S_{m,p}\equiv(2m+7)p^3\pmod{p^4}.
$$
It follows that
$$
 S_{m,p}\equiv0\pmod{p^4}
 \quad\Longleftrightarrow\quad
 p\mid2m+7.
$$
Together with \eqref{eq:num-equiv}, this proves the asserted
classification for every odd prime.

For $p=2$, the sum has one term:
$$
 S_{m,2}=\binom11^2\binom21^2=4.
$$
Thus $u_m(1)=4$, so $2^4\nmid u_m(1)$ and
$v_2(u_m(1))=2$.  This agrees with the criterion because
$2-1\mid2m+4$, whereas $2\nmid2m+7$.

It remains to record the valuation at an odd exceptional prime.
Equation \eqref{eq:harmonic-pairing} shows that $p\mid H_r$ in
$\Zloc$.  Equation \eqref{eq:first-reduction} therefore gives
$p^3\mid S_{m,p}$.  On the other hand, if
$p-1\mid2m+4$ and $p\nmid2m+7$, then \eqref{eq:key} and
\eqref{eq:power-sum} show that $S_{m,p}\not\equiv0\pmod{p^4}$.
Consequently
$$
 v_p\bigl(u_m(p-1)\bigr)=3.
$$

Finally, if $p\in\mathcal E_m$, then $p-1$ is a positive divisor of
$2m+4$.  Therefore
$$
 p-1\leq2m+4,\qquad p\leq2m+5.
$$
This proves the finiteness assertion and completes the proof of
Theorem~\ref{thm:main}.

\end{document}
